% English introduction replaced on 2026-11-06 (README (174)); verbatim, not compiled.
\section{Introduction}
\label{sec:introduction}

Positive-data grammatical inference asks when a language can be identified
in the limit from positive examples alone \cite{Gold1967,Angluin1980}.
Distributional learning treats words, or tuples of words, that occur in
common contexts as candidate constituents.  It began with substitutable
CFGs \cite{ClarkEyraud2007,Yoshinaka2008} and was extended to multiple
context-free grammars (MCFGs), whose tuple-valued nonterminals express
discontinuous and cross-serial dependencies
\cite{SekiEtAl1991,Yoshinaka2009,Yoshinaka2011,Kallmeyer2010}.
Yoshinaka's finite-signature framework \cite{Yoshinaka2015General} gives
positive-data identification with characteristic samples of polynomially
many words.  Polynomially many words may still be exponentially long, so we
measure a finite sample $K$ by its total encoded length
$\|K\|_+=\sum_{w\in K}(|w|+1)$, following the thick-data perspective of
\cite[Section~2.5.2]{EyraudHeinzYoshinaka2016}.  Qualitative identification
is inherited from Yoshinaka's framework; our subject is the length of
characteristic samples relative to an MCFG presentation of the target, and the
question is: \emph{for which comparison biases is this length polynomial, and
does the answer depend on the fan-out?}

\paragraph{The setting.}
Before any data are seen, we fix a fan-out bound $f$ and a surjective
finite-semigroup morphism $s:\Sigma^+\twoheadrightarrow\mathsf S$, a
\emph{typing}.  The reconstruction compares observed nonempty tuple components
only when their $s$-values agree, and the target class $\Cmcf{f}{s}$
consists of the $f$-MCFLs that are tuple-substitutable under this
restriction (Definition~\ref{def:class-cmcf}).  Already at fan-out two a
group typing makes this class strictly larger than Yoshinaka's untyped class,
with non-context-free witnesses on which the untyped reconstruction
over-generalizes (Proposition~\ref{prop:typed-separation}).  The typing is a
class-level bias in the sense of \cite{CosteTyping2004}, not chosen per
target; other biases restrict derivations by control sets
\cite{Takada1988,Takada1995} or restrict the evidence used to compare
factors \cite{CosteGaretNicolas2012,CosteNicolas2019}.  Typical typings are
prefix--suffix windows, values in a finite group, their products, and DFA
transition morphisms.  Group values compose in
the order of the blocks, so for noncommutative groups such as $S_3$ the order of
interleaved components matters.  Local groups, finite semigroups $\mathsf T$ with
$e\mathsf Te$ a group for every idempotent $e$, add group contributions of
sliding factors: over $\{\mathtt a,\mathtt b\}$, recording the parity of the
number of occurrences of $\mathtt{ab}$ is a local-group typing, but no
prefix--suffix window combined with a group morphism determines it: for each
window size and group morphism there are $m,n$ such that
$\mathtt b^m\mathtt a^{n}\mathtt b^m$ and $\mathtt b^{2m}$ agree on both but
differ in that parity
(Proposition~\ref{prop:li-g-strict}).  Windows, groups and their products are local
groups as well, but a DFA transition morphism need not be: if one letter acts as the
identity and another as a constant map on at least two states, the two transformations
form a copy of $U_1$, the obstruction of Theorem~\ref{thm:C} below.  For a reduced good presentation $G$ we
use its size $\|G\|$ and its rule thickness $\bar\theta_G$, the largest length
of a shortest sentence using a given rule.  Lower bounds are stated for every
set-driven operator with characteristic samples throughout the class; upper
bounds for Yoshinaka's reconstruction operator with typed links (Section~\ref{sec:learner}).

\paragraph{Main results.}
Fix the alphabet, the fan-out $f$, a rule-rank bound $r_{\max}\ge2$ and the
typing $s$.

\begin{mainthm}[Local groups]
\label{thm:A}
If the positive image of $s$ is a local group, every $L\in\Cmcf{f}{s}$ has
characteristic data of total length $O(\|G\|^2\bar\theta_G)$ for every reduced
good presentation $G$ \textup(Theorem~\ref{thm:local-group-polydata}\textup).
\end{mainthm}

\begin{mainthm}[The linear factor]
\label{thm:B}
\begin{enumerate}[label=\textup{(\roman*)},leftmargin=*]
\item If $\Sigma$ has at least $4f+1$ letters, or $f=1$ and $\Sigma$ has at least
two letters, then for every typing $s$ there are finite targets in
$\Cmcf{f}{s}$, with presentations within a constant factor of the smallest
reduced good ones, whose characteristic samples all have length
$\Omega(\|G\|\bar\theta_G)$
\textup(Proposition~\ref{prop:linear-factor-all-fanout}\textup).
\item For the trivial typing, Yoshinaka's bound $O(\|G\|\bar\theta_G)$ on characteristic
data \cite{Yoshinaka2011} is attained up to a constant factor: under the alphabet
condition of \textup{(i)}, by the targets of \textup{(i)}
\textup(Corollary~\ref{cor:trivial-exact}\textup).
\end{enumerate}
\end{mainthm}

\begin{mainthm}[The boundary at fan-out one]
\label{thm:C}
Let $f=1$ and $r_{\max}\ge2$.
\begin{enumerate}[label=\textup{(\roman*)},leftmargin=*]
\item A pseudovariety of finite semigroups admits uniform polynomial
characteristic-data length iff it is contained in $\mathbf{LG}$
\textup(Theorem~\ref{thm:fanout-one-data}\textup).
\item Outside $\mathbf{LG}$ the failure is exponential: for a typing $s_U$ onto
the two-element semilattice $U_1$ there are targets in $\Cmcf{1}{s_U}$ with
reduced good presentations of size and rule thickness $O(n)$ whose
characteristic samples all have length at least $3\cdot2^n-1$, for every
set-driven operator with characteristic samples throughout the class
\textup(Theorem~\ref{thm:fanout-one-data}\textup).
\end{enumerate}
\end{mainthm}

In the terminology of identification in polynomial time and thick data
\cite{Yoshinaka2008,EyraudHeinzYoshinaka2016}, Theorems~\ref{thm:A}
and~\ref{thm:C} say that for local-group typings the classes
$\Cmcf{f}{s}\cap\mathsf{MCFL}(f,r_{\max})$ are so identifiable at every fan-out,
and that at fan-out one $\mathbf{LG}$ is exactly the pseudovariety boundary of
this property (Corollary~\ref{cor:iptd}).  An MCFG of
fan-out one is a context-free grammar, so Theorem~\ref{thm:C} is a statement
about context-free languages; it turns the sufficient condition of the
companion manuscript \cite{kuriyamaCFG} into a characterization.

\paragraph{The remaining question.}
Theorem~\ref{thm:A} gives the \emph{if} direction of Theorem~\ref{thm:C}(i) at
every fan-out.  Whether $\mathbf{LG}$ is still the boundary at fan-out $f\ge2$
is open (Open Problem~\ref{op:fanout-two}).  The witness family of
Theorem~\ref{thm:C}(ii) does not decide it: by a two-hole transfer of powers,
under every fixed typing all but finitely many of its members, and their images
under letter-disjoint morphisms, leave the class at fan-out two, and no member
of the class keeps their long word as the only word of its type
(Theorem~\ref{thm:doubling-no-lift}).  A typing that needs superpolynomial data
at fan-out one but polynomial data at fan-out two would show that raising the
fan-out can make learning strictly cheaper.

\paragraph{Approach.}
Refining the grammar by types and applying untyped estimates to the
refinement can be expensive: for a fixed two-element typing there are
presentations $G_n$ of $\{\mathtt a,\mathtt c\}^+$ with $\|G_n\|=O(n)$ and
rule thickness at most $2$ whose full yield-type refinement has thickness at
least $2^n$ (Proposition~\ref{prop:typed-thickness-exponential}).  We
therefore bound characteristic data in terms of the original presentation,
and we use distribution-equivalent representatives rather than generated
witnesses (Corollary~\ref{cor:distribution-representatives}).
To obtain short tuples of prescribed type, we replace long subderivations
by short ones that need not preserve the group values of the replaced
subtrees; every regime is an instance of one shortening principle
(Lemma~\ref{lem:heavy-node-count}(ii)), and only the word-level replacement step
changes.  At fan-out one each replaced subtree occupies one contiguous block,
and a pigeonhole on prefix products of the group values selects a replaceable
interval.  In an MCFG derivation of fan-out $f\ge2$ a subtree contributes up to $f$
blocks to different root components, interleaved with blocks of other subtrees,
and for noncommutative groups the effect of simultaneous replacements depends on
this order.  The key step, with no fan-out-one counterpart, is an
\emph{owner-wise simultaneous selection} of replacements that preserves the
ordered group product of every root component at once
(Proposition~\ref{prop:block-replacement-interface}).  For
local-group typings the group value is a product over sliding factors, which
straddle block boundaries; selecting owners whose blocks lie far apart and
enlarging each block by unchanged boundary letters reduces this case to the
same selection (Lemma~\ref{lem:lg-buffered-owners}).  The lower bounds use
marker families all of whose subsets stay in the class
(Proposition~\ref{prop:linear-factor-all-fanout}) and rare types
(Proposition~\ref{prop:rare-type}).  A second ingredient without a fan-out-one
counterpart, a two-hole transfer of powers, excludes the exponential fan-out-one
obstruction at fan-out two (Lemma~\ref{lem:power-transfer},
Theorem~\ref{thm:doubling-no-lift}).

\paragraph{Relation to prior work.}
Yoshinaka bounded characteristic data of $(k,\ell)$-substitutable CFGs with
a thickness parameter \cite{Yoshinaka2008}, learned multidimensionally
substitutable MCFGs from positive data \cite{Yoshinaka2009,Yoshinaka2011},
and gave the finite-signature learning principle with finite witness
cardinality \cite{Yoshinaka2015General}; the qualitative results and the
Const/Link/Comp construction used here are his.  The finite-semigroup
structure theory and local-group representations are classical
\cite{Straubing1985,Tilson1987,AlmeidaEtAl2009}.  Our contribution is
the encoded-length analysis and the boundary results above.  Other
mildly context-sensitive formalisms are learned from positive data in
\cite{OatesEtAl2006}, and MCFGs with queries or dual context sets in
\cite{Yoshinaka2010,ClarkYoshinaka2014,KanazawaYoshinaka2023}; the MCFL
substitution lemma of \cite{DuncanEtAl2026Substitution} is a membership
condition, not a learning criterion.

\paragraph{Relation to the companion CFG manuscript.}
The companion manuscript \cite{kuriyamaCFG} treats fan-out one.  In short, it
proves polynomial data for typings in $\mathbf{LI}\vee\mathbf{Grp}$ as a
sufficient condition and exponential data for one typing onto $U_1$; here the
fan-out-one boundary is determined exactly ($\mathbf{LG}$), and the upper bound
for $\mathbf{LG}$ and the linear-factor lower bound hold at every fan-out.
Table~\ref{tab:companion-overlap} separates its results from ours (numbers in
the second column refer to the companion).  No result of the companion is used
as a premise: instead of its bracketed doubling family
\cite[Proposition~6.4]{kuriyamaCFG} we use a three-letter variant, proved in
Appendix~\ref{app:bracketed-family}, and its expressive results (counter and
Dyck obstructions, comparison with congruential grammars) are not used.

\begin{table}[htbp]
\caption{Results of the companion CFG manuscript \cite{kuriyamaCFG}
(fan-out one) and of the present paper.}
\label{tab:companion-overlap}
\centering
\footnotesize
\setlength{\tabcolsep}{3pt}
\renewcommand{\arraystretch}{1.15}
\begin{tabularx}{\linewidth}{@{}>{\raggedright\arraybackslash}p{0.27\linewidth}>{\raggedright\arraybackslash}p{0.25\linewidth}>{\raggedright\arraybackslash}X@{}}
\hline
Topic & Companion \cite{kuriyamaCFG} & Here\\
\hline
Typed reconstruction, Gold learner, typed-thickness witness bound
 & Theorem~4.2, Theorems~5.2--5.3, Corollary~5.4; Lemma~6.2, Corollary~6.3
 & Tuple versions (Section~\ref{sec:learner},
   Theorem~\ref{thm:data-from-typed-thickness}); qualitative results credited to
   Yoshinaka; not new\\
Window and window--group bounds
 & Proposition~3.5, Lemma~7.1, Theorem~7.2, Lemma~7.4, Corollary~7.5
 & Not new at $f=1$; new for $f\ge2$: owner-wise replacement across interleaved
   components (Proposition~\ref{prop:block-replacement-interface})\\
Local-group typings
 & Not treated (only $\mathbf{LI}\vee\mathbf{Grp}$; Lemma~7.6)
 & New for data length, at every $f$ (Theorems~\ref{thm:lg-characterization}
   and~\ref{thm:local-group-polydata}); the characterization of $\mathbf{LG}$ by
   $U_1$ and the semigroup algebra are classical \cite{AlmeidaEtAl2009}\\
Doubling family under a typing onto $U_1$
 & Proposition~6.4, Corollary~6.6, Appendix~C (five letters, brackets)
 & Three-letter variant with the same bounds, proved here
   (Proposition~\ref{prop:three-letter-family})\\
Rare-type lower bound
 & Proposition~6.4(iii)--(iv), for a single word
 & Every realized type and every $f$ (Proposition~\ref{prop:rare-type})\\
Fan-out-one pseudovariety boundary
 & Sufficiency for $\mathbf{LI}\vee\mathbf{Grp}$ only (Remark~7.7)
 & New (Theorem~\ref{thm:fanout-one-data})\\
Typed versus untyped class
 & Separations at $f=1$
 & At $f\ge2$, a non-context-free member outside Yoshinaka's class
   (Proposition~\ref{prop:typed-separation})\\
Necessity of the factor $\|G\|\bar\theta_G$
 & ---
 & New at every $f$ (Proposition~\ref{prop:linear-factor-all-fanout})\\
Fan-out two: power transfer
 & ---
 & New (Theorem~\ref{thm:doubling-no-lift})\\
\hline
\end{tabularx}
\end{table}

\paragraph{Scope of the claims.}
All bounds concern the total encoded length of characteristic samples for
the specified set-driven operator, relative to the stated presentation;
they bound neither an intrinsic grammar size nor a stage of Gold
convergence, and their exponents are discussed at the end of
Section~\ref{sec:conclusion}.  Lower bounds on data come only from
distinguishing targets (Propositions~\ref{prop:rare-type}
and~\ref{prop:linear-factor-all-fanout}); thick refinements
(Proposition~\ref{prop:typed-thickness-exponential}) limit the witness
construction, not the data (Corollary~\ref{cor:distribution-representatives}).
The typing is fixed before the target is seen; refining it enlarges the class
but changes the reconstruction operator, so we neither infer nor optimize it.
Table~\ref{tab:scope} lists, for each main claim, the parameters that are
fixed, the operators it concerns and its side conditions; constants in upper
bounds depend on the fixed parameters only.

\begin{table}[htbp]
\caption{Scope of the main claims.  ``Canonical'' means both
$\mathcal B^{\le r_{\max}}_{f,s}$ and $\mathcal B_{f,s}$ (with one common sample);
``class-wide'' means every set-driven operator with characteristic samples
throughout the class.  $G$ is a reduced good presentation of $L\cap\Sigma^+$ of
fan-out $\le f$ and rank $\le r_{\max}$.}
\label{tab:scope}
\centering
\footnotesize
\setlength{\tabcolsep}{3pt}
\renewcommand{\arraystretch}{1.15}
\begin{tabularx}{\linewidth}{@{}>{\raggedright\arraybackslash}p{0.2\linewidth}>{\raggedright\arraybackslash}p{0.22\linewidth}>{\raggedright\arraybackslash}p{0.17\linewidth}>{\raggedright\arraybackslash}X@{}}
\hline
Claim & Fixed & Operators & Bound and conditions\\
\hline
Theorem~\ref{thm:A} (Theorem~\ref{thm:local-group-polydata})
 & $f$, $r_{\max}$, $\Sigma$, $s$ with $s(\Sigma^+)\in\mathbf{LG}$
 & canonical
 & upper $O(\|G\|^2\bar\theta_G)$, every $L\in\Cmcf{f}{s}$ and every $G$\\
Theorem~\ref{thm:B}(i) (Proposition~\ref{prop:linear-factor-all-fanout})
 & $f$, any $s$
 & canonical ($r\ge1$), class-wide
 & lower $\Omega(\|G\|\bar\theta_G)$ for finite targets; needs $|\Sigma|\ge4f+1$ ($|\Sigma|\ge2$ if $f=1$);
   presentations of rank two\\
Theorem~\ref{thm:B}(ii) (Corollary~\ref{cor:trivial-exact})
 & $f$, $r_{\max}\ge1$, $s_{\mathrm{triv}}$
 & $\mathcal B_{f,s_{\mathrm{triv}}}$
 & upper $O(\|G\|\bar\theta_G)$ (Yoshinaka \cite{Yoshinaka2011}); attained if $r_{\max}\ge2$ and $|\Sigma|\ge4f+1$ ($\ge2$ if $f=1$)\\
Theorem~\ref{thm:C}(i) (Theorem~\ref{thm:fanout-one-data})
 & $f=1$, $r_{\max}\ge2$, a pseudovariety
 & $\mathcal B^{\le r_{\max}}_{1,s}$ (Definition~\ref{def:uniform-polydata})
 & uniform polynomial data iff $\mathbf V\subseteq\mathbf{LG}$\\
Theorem~\ref{thm:C}(ii)
  & $f=1$, $s_U$ on three letters
 & canonical, class-wide
 & lower $3\cdot2^n-1$ for presentations of size and thickness $O(n)$\\
Corollary~\ref{cor:iptd}
 & $f$, $r_{\max}$, $s$
 & set-driven
 & IPTtD for $\mathbf{LG}$ at every $f$; at $f=1$ exactly for $\mathbf{V}\subseteq\mathbf{LG}$\\
Theorem~\ref{thm:doubling-no-lift}
 & any $f\ge2$, any $s$
 & none (membership)
  & the family of Theorem~\ref{thm:C}(ii) and its letter-disjoint images leave $\Cmcf{f}{s}$, and their forced word is never the only one of its type there; no data bound at $f\ge2$\\
Open
 & 
 & 
 & whether $\mathbf{LG}$ is the boundary at $f\ge2$ (Open Problem~\ref{op:fanout-two})\\
\hline
\end{tabularx}
\end{table}

\paragraph{Organization.}
Section~\ref{sec:prelim} fixes the setting and Section~\ref{sec:learner}
recalls the reconstruction operator and its qualitative guarantees, and
separates the typed from the untyped class (Section~\ref{subsec:typed-separation}).
Section~\ref{sec:char-data} reduces characteristic data to typed thickness and
to distribution representatives.  Section~\ref{sec:upper} proves
Theorem~\ref{thm:A}, Section~\ref{sec:linear} proves Theorem~\ref{thm:B}, and
Section~\ref{sec:fanout-one} proves Theorem~\ref{thm:C} and restates the
boundary in terms of identification in polynomial time and thick data.
Section~\ref{sec:fanout-two} shows that the witness family does not carry over
to fan-out two, and Section~\ref{sec:conclusion} discusses the exponents and the
open problem.  The appendices contain the technical proofs.

